\section{Experiments}
\label{sec:exp}

\subsection{Experimental Setup}
\label{sec:exp:setup}

\textbf{Topology and Traffic.} The main experiments are conducted on the GEANT backbone topology ($22$ nodes, $36$ links). To train a policy capable of generalizing across different network states, all link attributes are randomly regenerated at the start of each episode: bandwidth is sampled from $\text{Uniform}[30,100]$ and delay from $\text{Uniform}[1,20]$. Traffic generation uses QoS intent $\phi$ as the dominant variable: $\phi$ is generated from a flipped Pareto distribution ($\alpha\!=\!0.2$), so that the overwhelming majority of flows fall in the low-$\phi$ region (mouse flows with strict bandwidth and delay constraints) and only a small number are high-$\phi$ flows (elephant flows with relaxed constraints). Bandwidth demand $b$ is derived from $\phi$ via a positively correlated linear mapping (Eq.~\eqref{eq:bw_gen}), with $b_{\text{max}}\!=\!40$ and Gaussian noise standard deviation equal to $0.15$ times the span, so that $\phi$ and $b$ remain strongly positively correlated (Pearson $r\!\approx\!0.90$) but not fully deterministic -- $\phi$ still carries independent information beyond $b$.

\textbf{Sequence Length.} Each episode contains $W=150$ flows arriving sequentially, requiring the policy to maintain service quality under long-range credit assignment. Randomized topologies and flow generation ensure that the agent cannot memorize fixed paths and must make differentiated routing decisions based on current link states and the intent of each flow.

\textbf{Baselines.} We compare against four classical routing strategies: OSPF (shortest hop count), ECMP (equal-cost multipath random), CSPF (online constrained shortest path), and RANDOM. All methods are evaluated on the same $200$ test traffic sequences (episodes) and averaged; the PPO policy uses deterministic actions (taking the highest-probability candidate among the $K$ paths, without exploration sampling) to ensure reproducibility.

\textbf{Evaluation Metrics.} Mean Reward (Mean $R$), Bandwidth Satisfaction Score (BW Score, i.e., the mean of per-flow bandwidth satisfaction ratios $\rho_t$), and Average End-to-End Delay (Avg Delay, i.e., the mean of effective delays $\hat D_t$ (Eq.~\eqref{eq:delay}) across all flows). Bandwidth scores are further broken down by flow category (small mouse / mid / big elephant) to examine whether the policy differentiates treatment by QoS intent.

\subsection{Hyperparameters and Implementation Details}
\label{sec:exp:hyperparam}

All hyperparameters denoted symbolically in Sections~\ref{sec:problem}--\ref{sec:method} are assigned specific values in Table~\ref{tab:hyperparam}, with the rationale as follows. For model architecture, the hidden dimension $h=256$ and 4-layer GNN follow the common configuration range of $h\in[64,256]$ and $L\in[2,5]$ seen in RouteNet~\cite{Rusek2020_JSAC_RouteNet} and GNN-FiLM~\cite{Brockschmidt2020_GNNFiLM} in routing scenarios; the single attention head matches the default setting of prior work~\cite{Shi2021_UniMP}. For PPO training, $\gamma=1.0$ is used because the task involves finite-horizon episodes (episodic); the remaining clipping coefficient, GAE parameters, and value coefficient follow the defaults of the original PPO papers~\cite{Schulman2017_PPO,Schulman2016_GAE}; the entropy coefficient $c_e=0.01$ and minibatch size were determined via grid search. For the task, Pareto $\alpha=0.2$ makes the proportion of low-$\phi$ flows exceed 80\%, matching the heavy-tailed characteristics of real traffic reported by measurement studies~\cite{CISCO2023_GlobalTraffic}; the ratio of the flow bandwidth upper bound $b_{\text{max}}=40$ to the link capacity lower bound $30$ means a single flow can create a bottleneck, providing a learning signal for the policy to reserve capacity; the delay normalization constant $\tau=38.3$ is taken as the P25 delay of the GEANT topology. Model architecture and training coefficients are identical across all variants; all variants share the same hyperparameters; link attributes are randomly resampled per episode. Implementation is based on PyTorch, TorchRL, and PyTorch Geometric.

\begin{table}[t]
\centering
\caption{Hyperparameter summary. All variants share these values.}
\label{tab:hyperparam}
\setlength{\tabcolsep}{3pt}
\footnotesize
\begin{tabular}{p{1.6cm}p{3.6cm}p{2.8cm}}
\toprule
Category & Hyperparameter & Value \\
\midrule
\multirow{8}{*}{Model Arch.}
 & Hidden dim $h$              & $256$ \\
 & GNN layers $L$            & $4$ \\
 & Attention heads             & $1$ \\
 & Candidate paths $K$          & $16$ \\
 & $\mathrm{MLP}_{\text{lift}}$ hidden & $[64,128]$, GELU \\
 & Scorer feature MLP         & $[512,512]$, ELU$+$LN \\
 & Scoring head MLP             & $[256,128,64]$, ELU$+$LN \\
 & PathPooler            & Single-layer LSTM \\
\midrule
\multirow{8}{*}{PPO Train.}
 & Discount $\gamma$           & $1.0$ \\
 & GAE $\lambda$          & $0.95$ \\
 & Clipping coeff.\ $\epsilon$     & $0.2$ \\
 & Value coeff.\ $c_v$          & $0.5$ \\
 & Entropy coeff.\ $c_e$            & $0.01$ \\
 & Batch frames $B$            & $3000$ \\
 & PPO epochs $E$ / minibatch & $5$ / $512$ \\
 & Learning rate $\eta$ / optimizer   & $3\times10^{-4}$ / Adam \\
\midrule
\multirow{11}{*}{Task}
 & Topology                  & GEANT ($22$ / $36$) \\
 & Link BW gen.           & $\text{Uniform}[30,100]$ \\
 & Link delay gen.           & $\text{Uniform}[1,20]$ \\
 & Sequence length $W$            & $150$ \\
 & Flow BW upper $b_{\text{max}}$ & $40$ \\
 & Flow BW lower $b_{\text{floor}}$ & $0$ \\
 & $\phi$ dist. Pareto $\alpha$ & $0.2$ \\
 & BW noise $\sigma$      & 6.0 \\
 & Delay norm.\ $\tau$  & $38.3$ (GEANT P25)\\
 & Training frames $T$         & $2\times10^{6}$ \\
 & Gradient clip           & $1.0$ \\
\bottomrule
\end{tabular}
\end{table}

\subsection{Experiment I: Flow-Aware Representations Significantly Outperform Classical Routing Algorithms}
\label{sec:exp:main}

We first compare QoS-FiLM against classical routing algorithms to quantify the advantage of flow-aware representations. To isolate the ``learned policy vs.\ classical algorithm'' comparison, this section contrasts QoS-FiLM against the four classical baselines; ablation comparisons with the noFiLM-Tf and MLP (no message-passing) variants are deferred to Section~\ref{sec:exp:repr}.

At $W\!=\!150$, \textcolor{red}{FiLM-GNN} outperforms all classical baselines on all three metrics of mean reward, bandwidth satisfaction, and delay (Fig.~\ref{fig:main_bars}): its mean reward of $+0.121$ surpasses the strongest classical baseline CSPF ($-0.046$) by $+0.167$, BW Score of $0.759$ exceeds CSPF ($0.473$) by $0.286$, and average delay drops from $130.0$ to $112.9$~ms. Although CSPF has the lowest delay ($125.9$), it does so at the sacrifice of bandwidth satisfaction (BW Score only $0.473$), resulting in the lowest overall reward ($-0.208$)\textcolor{red}{这个值怎么跟本段前面的值不一致？} — looking at delay alone cannot reflect the global trade-off in service quality. This advantage holds across all three topologies (Fig.~\ref{fig:main_bars}): QoS-FiLM outperforms the respective strongest baseline CSPF on GEANT ($+0.121$ vs.\ $-0.046$), NOBEL-EU ($-0.108$ vs.\ $-0.183$), and ABILENE ($+0.122$ vs.\ $+0.041$), demonstrating that the advantage of flow-aware representations is robust across topologies.

\textbf{Source of Advantage: Non-Myopic Long-Range Resource Allocation.} Classical algorithms (including the ``stepwise optimal'' CSPF) make independent greedy routing decisions for each flow and lack awareness of reserving capacity for subsequent flows; under the long sequence length of $W\!=\!150$, early-arriving flows deplete critical link capacity, leaving later flows with no viable path. Bucketing by $\phi$ makes this clear: in the low-intent bucket ($\phi<0.2$, the most strictly constrained mouse flows), \textcolor{red}{FiLM-GNN} achieves a reward of $+0.315$, far surpassing CSPF's $+0.070$; in the mid-intent bucket the two converge; in the high-intent bucket all methods score low — an inherent resource scarcity characteristic of the task. This confirms the argument of Section~\ref{sec:problem:why_rl}: routing is a sequential decision problem with long-range credit assignment, and CSPF, as the representative of ``stepwise optimal,'' is systematically surpassed by the learned non-myopic policy.

\begin{figure*}[t]
\centering
\subimg{p9b_bars_geant_linear_reward.pdf}\hfill
\subimg{p9b_bars_nobel_eu_reward.pdf}\hfill
\subimg{p9b_bars_abilene_reward.pdf}\\[4pt]
\subimg{p9b_bars_geant_linear_bw.pdf}\hfill
\subimg{p9b_bars_nobel_eu_bw.pdf}\hfill
\subimg{p9b_bars_abilene_bw.pdf}\\[4pt]
\subimg{p9b_bars_geant_linear_delay.pdf}\hfill
\subimg{p9b_bars_nobel_eu_delay.pdf}\hfill
\subimg{p9b_bars_abilene_delay.pdf}
\caption{Experiment I: Final performance bar chart comparison across three topologies. (a)--(c) mean reward (mean of 3 seeds, error bars $\pm$ std.\ dev.); (d)--(f) bandwidth service rate; (g)--(i) average delay (ms). Columns from left to right: GEANT ($W\!=\!150$), NOBEL-EU ($W\!=\!150$), ABILENE ($W\!=\!50$). Solid colored bars represent learned network variants (QoS-FiLM highlighted), gray hatched bars represent static baselines.}
\label{fig:main_bars}
\end{figure*}

\subsection{Experiment II: Representational Capacity Ablation -- Message Passing and Intent Conditioning}
\label{sec:exp:repr}

To isolate the contribution of each architectural component, we design three variants along the progressive path of ``no graph structure $\to$ with graph structure but no flow awareness $\to$ with graph structure and flow awareness'' (Fig.~\ref{fig:main_bars}), all sharing the same TransformerConv backbone and hyperparameters, differing only in their representational mechanism:
\begin{itemize}
\item \textbf{MLP} (no message passing): removes message passing, replacing $\mathrm{Conv}$ in Eq.~\eqref{eq:conv} with a per-node MLP (still retaining FiLM modulation), ignoring the edge set $E$ and link features $f$, to test whether graph structure information is necessary.
\item \textbf{noFiLM-Tf} (with message passing but intent only concatenated at the end): removes the $\gamma,\beta$ modulation of Eq.~\eqref{eq:film} (equivalent to $\gamma\equiv 0,\beta\equiv 0$); the intent vector $m_t$ does not enter message passing and is only concatenated with path embeddings at the final scorer, so all flows see the same graph.
\item \textbf{QoS-FiLM} (per-layer intent injection): our method, modulating node representations at every layer via the FiLM affine transformation of Eq.~\eqref{eq:film}.
\end{itemize}

\textbf{Message passing is foundational.} The MLP baseline scores worst on all three metrics (mean reward $-0.177$, BW Score $0.520$), with all flows degenerating into selecting the same path slot (Table~\ref{tab:film_w150_bw} shows converging scores across categories, approximating OSPF-equivalent behavior). Without neighbor aggregation, the model cannot perceive the global distribution of link utilization and cannot differentiate routing for flows with different intents -- confirming that graph structure information is indispensable for routing learning.

\textbf{Per-layer intent injection (FiLM) brings significant gains.} Under the same GNN backbone, per-layer intent-injected QoS-FiLM ($+0.121$) significantly outperforms end-only intent-concatenated noFiLM-Tf ($-0.098$), $\Delta R\!=\!+0.219$; BW Score improves by $+0.158$ ($0.759$ vs.\ $0.601$), and average delay drops from $157.8$ to $112.9$ (a reduction of $44.9$~ms). Fig.~\ref{fig:repr_training} shows the Eval Reward training curves of the two variants. The disadvantage of noFiLM-Tf stems from intent information being concatenated only at the end, unable to apply $\gamma\!\cdot\! x+\beta$ modulation to per-layer node representations, and thus unable to reshape the message-passing process according to flow intent. This trend is consistent across all three topologies (Fig.~\ref{fig:repr_training}): QoS-FiLM clearly leads noFiLM-Tf and MLP, and converges faster with smaller variance (NOBEL-EU: QoS-FiLM $-0.108$ vs.\ noFiLM-Tf $-0.282$; ABILENE: $+0.122$ vs.\ $+0.008$).

\begin{figure*}[t]
\centering
\subimg{p9b_groupA_geant_linear_eval.pdf}\hfill
\subimg{p9b_groupA_nobel_eu_eval.pdf}\hfill
\subimg{p9b_groupA_abilene_eval.pdf}\\[4pt]
\subimg{p9b_groupA_geant_linear_bw.pdf}\hfill
\subimg{p9b_groupA_nobel_eu_bw.pdf}\hfill
\subimg{p9b_groupA_abilene_bw.pdf}
\caption{Experiment II: Training curves of QoS-FiLM, noFiLM-Tf, and MLP on three topologies (mean of 3 seeds $\pm$ std.\ dev.\ confidence band). (a)--(c) Evaluation Reward; (d)--(f) Bandwidth Service Rate. Columns from left to right: GEANT ($W\!=\!150$), NOBEL-EU ($W\!=\!150$), ABILENE ($W\!=\!50$). Gray dashed lines indicate the final evaluation performance of classical baselines.}
\label{fig:repr_training}
\end{figure*}

Viewing by $\phi$ bucket (Table~\ref{tab:film_w150_phi}), QoS-FiLM achieves a delay of only $77.6$ in the latency-sensitive bucket ($\phi<0.2$), a reduction of $68.4$~ms compared to noFiLM-Tf ($146.0$); whereas in the latency-insensitive bucket ($\phi\geq 0.8$), QoS-FiLM's delay ($290.0$) is actually higher than noFiLM-Tf's ($225.9$). This ``low-down, high-up'' pattern is the expected behavior of FiLM conditioning: the policy discriminates among different $\phi$ values and actively transfers delay resources from high-intent flows to low-intent ones, achieving a significant improvement in service quality for the low-intent bucket at minimal cost. By bandwidth category (Table~\ref{tab:film_w150_bw}), QoS-FiLM's BW Score on small flows reaches $0.932$ (noFiLM-Tf only $0.669$, gain $+0.263$), and similarly leads substantially on mid flows ($0.857$ vs.\ $0.614$); on big flows QoS-FiLM is slightly lower ($0.348$ vs.\ $0.483$), but achieves ultra-high scores on small/mid flows at a relatively small cost to big flows, a favorable overall trade-off.

\begin{table}[t]
\centering
\caption{Experiment II ablation -- $\phi$-bucketed delay ($W\!=\!150$, 200 ep $\times$ 3 seeds)\protect\footnotemark[1].}
\label{tab:film_w150_phi}
\setlength{\tabcolsep}{3pt}
\begin{tabular}{lccc}
\toprule
Model & $\phi<0.2$ & $0.2\!\le\!\phi\!<\!0.8$ & $\phi\!\ge\!0.8$ \\
\midrule
\textbf{QoS-FiLM} & $\mathbf{77.6}$ & $190.0$ & $290.0$ \\
noFiLM-Tf        & $146.0$ & $181.7$ & $225.9$ \\
MLP              & $156.3$ & $183.7$ & $220.2$ \\
\bottomrule
\end{tabular}
\end{table}

\begin{table}[t]
\centering
\caption{Experiment II ablation -- BW categories ($W\!=\!150$, 200 ep $\times$ 3 seeds)\protect\footnotemark[1].}
\label{tab:film_w150_bw}
\setlength{\tabcolsep}{3pt}
\begin{tabular}{lccc}
\toprule
Model & $b<8$ & $8\!\le\!b\!<\!32$ & $b\!\ge\!32$ \\
\midrule
\textbf{QoS-FiLM} & $\mathbf{0.932}$ & $\mathbf{0.857}$ & $0.348$ \\
noFiLM-Tf        & $0.669$ & $0.614$ & $0.483$ \\
MLP              & $0.568$ & $0.531$ & $0.431$ \\
\bottomrule
\end{tabular}
\end{table}
\footnotetext[1]{All metrics are the mean summary of 3 random seeds on GEANT (200 ep $\times$ 3 seeds, $W\!=\!150$).\textcolor{red}{不要脚注吧？}}

\subsection{Experiment III: Graph Attention Operator Selection -- TransformerConv vs.\ GATv2Conv}
\label{sec:exp:conv}

We next compare two graph attention operators under the same FiLM framework. TransformerConv uses dot-product attention, while GATv2Conv uses dynamic additive attention~\cite{Brody2022_GATv2}, representing two different attention paradigms. Fig.~\ref{fig:conv_training} shows the training curve comparison of QoS-FiLM under the two operators (along with noFiLM-GATv2 as a reference) across three topologies; QoS-FiLM stably leads FiLM-GATv2 on all three topologies\textcolor{red}{这句话是不是有歧义？}.

\begin{figure*}[t]
\centering
\subimg{p9b_groupB_geant_linear_eval.pdf}\hfill
\subimg{p9b_groupB_nobel_eu_eval.pdf}\hfill
\subimg{p9b_groupB_abilene_eval.pdf}\\[4pt]
\subimg{p9b_groupB_geant_linear_bw.pdf}\hfill
\subimg{p9b_groupB_nobel_eu_bw.pdf}\hfill
\subimg{p9b_groupB_abilene_bw.pdf}
\caption{Experiment III: Training curves of QoS-FiLM, FiLM-GATv2, and noFiLM-GATv2 on three topologies (mean of 3 seeds $\pm$ std.\ dev.\ confidence band). (a)--(c) Evaluation Reward; (d)--(f) Bandwidth Service Rate. Columns from left to right: GEANT ($W\!=\!150$), NOBEL-EU ($W\!=\!150$), ABILENE ($W\!=\!50$).}
\label{fig:conv_training}
\end{figure*}

In terms of final policy quality (Fig.~\ref{fig:conv_training}), QoS-FiLM (TransformerConv$+$FiLM) is optimal in mean reward ($+0.121$) and bandwidth score ($0.759$), substantially leading FiLM-GATv2 ($-0.097$, $0.596$), with a BW Score gap of $+0.163$. The root cause of the performance difference between the two operators lies in their fundamentally different treatment of link attributes (delay, capacity, residual bandwidth $f_e=[d_e,c_e,b_e^{\text{res}}]$). TransformerConv injects link attributes into both the attention key and the message value: attention weight $\alpha_{i,j}=\mathrm{softmax}((\mathbf{W}_3\mathbf{x}_i)^{\top}(\mathbf{W}_4\mathbf{x}_j+\mathbf{W}_6\mathbf{e}_{ij})/\sqrt{d})$, message $\mathbf{W}_2\mathbf{x}_j+\mathbf{W}_6\mathbf{e}_{ij}$. The link feature $\mathbf{e}_{ij}$ independently modulates attention via multiplicative interaction (inner product) -- a high-residual-bandwidth link can receive higher attention weight without depending on target node features; simultaneously, the link feature directly constitutes the message content, so aggregated information naturally carries link quality signals. In contrast, GATv2Conv~\cite{Brody2022_GATv2} mixes link and node features in an additive manner: $\alpha_{i,j}\propto\exp(\mathbf{a}^{\top}\mathrm{LeakyReLU}(\mathbf{\Theta}_s\mathbf{x}_i+\mathbf{\Theta}_t\mathbf{x}_j+\mathbf{\Theta}_e\mathbf{e}_{ij}))$. Link features are summed with source and target node features before jointly entering the nonlinear activation, causing the link quality signal to be strongly coupled with node features -- the attention weight of a high-bandwidth link depends heavily on the neighboring node representations, diluting the independent discriminative power of link attributes. In a routing scenario, link residual bandwidth is a hard constraint signal for path selection and should independently and directly influence neighbor aggregation; TransformerConv's dual-channel injection of link attributes (key + value) naturally suits this need, whereas GATv2Conv's additive mixing weakens the independent expressiveness of link quality.

By $\phi$ bucket (Table~\ref{tab:conv_w150_phi}), QoS-FiLM similarly exhibits the ``low-down, high-up'' pattern described in Section~\ref{sec:exp:repr}: delay in the low-intent bucket is only $77.6$ (GATv2 is $135.2$, a reduction of $57.6$~ms), while delay in the high-intent bucket of $290.0$ is slightly higher than GATv2's $253.5$ -- TransformerConv identifies the $\phi$ semantics of flows more accurately than GATv2, transferring delay resources from high-intent to low-intent flows. By bandwidth category (Table~\ref{tab:conv_w150_bw}), QoS-FiLM's small/mid bandwidth scores ($0.932$/$0.857$) both substantially lead FiLM-GATv2 ($0.667$/$0.618$), with slightly lower big-flow scores ($0.348$ vs.\ $0.456$), achieving ultra-high small/mid flow scores at a relatively small cost to big flows.

\begin{table}[t]
\centering
\caption{Experiment III ablation -- $\phi$-bucketed delay ($W\!=\!150$, 200 ep $\times$ 3 seeds).}
\label{tab:conv_w150_phi}
\setlength{\tabcolsep}{3pt}
\begin{tabular}{lccc}
\toprule
Model & $\phi<0.2$ & $0.2\!\le\!\phi\!<\!0.8$ & $\phi\!\ge\!0.8$ \\
\midrule
\textbf{QoS-FiLM} & $\mathbf{77.6}$ & $190.0$ & $290.0$ \\
FiLM-GATv2       & $135.2$ & $189.1$ & $253.5$ \\
MLP              & $156.3$ & $183.7$ & $220.2$ \\
\bottomrule
\end{tabular}
\end{table}

\begin{table}[t]
\centering
\caption{Experiment III ablation -- BW categories ($W\!=\!150$, 200 ep $\times$ 3 seeds).}
\label{tab:conv_w150_bw}
\setlength{\tabcolsep}{3pt}
\begin{tabular}{lccc}
\toprule
Model & Small BW & Mid BW & Big BW \\
\midrule
\textbf{QoS-FiLM} & $\mathbf{0.932}$ & $\mathbf{0.857}$ & $0.348$ \\
FiLM-GATv2       & $0.667$ & $0.618$ & $0.456$ \\
MLP              & $0.568$ & $0.531$ & $0.431$ \\
\bottomrule
\end{tabular}
\end{table}

\subsubsection{Strong Synergy Between FiLM and Operator: Two-Dimensional Non-Additivity}
\label{sec:exp:synergy}

We further examine whether operator choice and conditioning interact synergistically. We supplement the noFiLM-GATv2 variant, forming a $2\times2$ comparison table of operator $\times$ conditioning (Table~\ref{tab:synergy_2x2})\textcolor{red}{（可以不改，思考一下）MLP有FiLM，但是这个表里没有MLP的结果；是不是能跟noFilM-Tf对比说明图结构和Film哪个更重要？}.

The key finding is: \emph{the gain from FiLM strongly depends on the underlying operator}. On TransformerConv, introducing FiLM brings a massive improvement of $\Delta R\!=\!+0.219$ ($-0.098\!\to\!+0.121$); but on GATv2, the gain from FiLM is merely $\Delta R\!=\!+0.007$ ($-0.104\!\to\!-0.097$), nearly ineffective. This indicates that the two dimensions are not independently additive but exhibit strong synergy: the aforementioned difference in how link attributes are handled amplifies the gap in FiLM effectiveness -- TransformerConv independently injects link attributes into attention keys and message values, providing a high-quality, link-level differentiable representational foundation for FiLM's per-layer intent modulation; whereas GATv2Conv mixes link attributes into node features, so even if FiLM accurately injects flow intent at each layer, the intent signal struggles to be converted into effective path differentiation through the diluted link quality representations. Therefore, the combination of TransformerConv and FiLM is not a simple superposition of two independent advantages, but a synergistic product in which neither can be omitted.

\begin{table}[t]
\centering
\caption{Operator $\times$ conditioning $2\times2$ comparison table (Mean $R$, $W\!=\!150$, 200 ep $\times$ 3 seeds). Values in parentheses show the FiLM gain $\Delta R$ on the same operator.}
\label{tab:synergy_2x2}
\setlength{\tabcolsep}{6pt}
\begin{tabular}{lcc}
\toprule
 & TransformerConv & GATv2Conv \\
\midrule
w FiLM & $\mathbf{+0.121}$ & $-0.097$ \\
w/o FiLM & $-0.098$ & $-0.104$ \\
\midrule
$\Delta R$ (FiLM gain) & $\mathbf{+0.219}$ & $+0.007$ \\
\bottomrule
\end{tabular}
\end{table}

The parameter overhead of FiLM is an operator-independent constant of $+6{,}144$ (relative overhead only $+0.6\%\sim+2.0\%$, coming from two intent$\to$hidden-dim linear layers $\mathrm{MLP}_\gamma,\mathrm{MLP}_\beta$). On the TransformerConv backbone, FiLM trades less than $1\%$ of parameters for an improvement of $\Delta R\!=\!+0.219$, yielding high cost-effectiveness.

\subsection{Experiment IV: Interpretability -- How Flow Intent Influences Representation and Decision}
\label{sec:exp:interp}

To further understand the mechanism of FiLM conditioning, this section examines the essential differences between FiLM and noFiLM-Tf in responding to flow intent from two perspectives: the \emph{representation level} and the \emph{decision level}.

\textbf{Representation level: intent-driven hidden feature sensitivity.} Fixing the network topology and link utilization state, we inject flows with gradually varying QoS intent $\phi$ into the same (src, dst) node pair and observe the hidden features produced by the encoder. We compute intent sensitivity (total change in hidden representations caused by varying $\phi$ under the same state) over $8$ random network states and $\phi\!\in\!\{0.05,0.2,0.5,0.8,0.95\}$: the FiLM encoder has a sensitivity as high as $46.46$, while noFiLM-Tf is only $1.4\times10^{-6}$ (numerically zero), a difference of seven orders of magnitude (Table~\ref{tab:interp_repr}). Furthermore, as $\phi$ gradually increases from the baseline ($0.05\!\to\!0.95$), the drift of FiLM representations monotonically increases ($0\!\to\!8.6\!\to\!25.8\!\to\!40.5\!\to\!46.5$), indicating that intent is smoothly and continuously encoded into node representations; whereas the drift of noFiLM-Tf remains zero throughout. This confirms the ablation design at the architectural level (noFiLM indeed cannot perceive intent) and provides direct quantitative evidence that FiLM injects per-flow intent into node representations. Fig.~\ref{fig:embedding_separation} further shows the PCA projection of hidden features from both encoders: FiLM's representations migrate continuously with $\phi$, while noFiLM-Tf's representations are completely overlapping.

\begin{table}[t]
\centering
\caption{Experiment IV representation level: hidden representation change when only $\phi$ varies under the same network state (GEANT, mean over $8$ states).}
\label{tab:interp_repr}
\setlength{\tabcolsep}{5pt}
\begin{tabular}{lcc}
\toprule
 & Intent Sensitivity & Drift ($\phi\!:0.05\!\to\!0.95$) \\
\midrule
\textbf{QoS-FiLM}  & $\mathbf{46.46}$ & $0\to8.6\to25.8\to40.5\to46.5$ \\
noFiLM-Tf         & $1.4\times10^{-6}$ & All zero \\
\bottomrule
\end{tabular}
\end{table}

\begin{figure*}[t]
\centering
\subimgw{p9b_hidden_diff_film_tf.pdf}\hfill
\subimgw{p9b_hidden_diff_nofilm_tf.pdf}
\caption{Comparison of hidden features for an extreme small flow and an extreme large flow under the same network state and (src, dst) pair.\textcolor{red}{后面的这个子标题直接放到上面的(a)(b)后面} (a) QoS-FiLM; (b) noFiLM-Tf.}
\label{fig:embedding_separation}
\end{figure*}

\textbf{Decision level: intent-driven path selection diversity.} Differences at the representation level ultimately manifest as changes in routing behavior. We compute the entropy of the path-slot selection distribution (in bits; higher indicates more diverse selection, more differentiation by state and intent) for both variants over $200$ test episodes comprising $460$ (src, dst) pairs. FiLM's mean entropy is $2.235$, significantly higher than noFiLM-Tf's $1.962$ (paired $t\!\approx\!14.7$, highly significant\textcolor{red}{这里有t-test，前面V-C里面是不是最好也要有}); on $76\%$ of (src, dst) pairs, FiLM's entropy is higher, and FiLM on average uses $10.4$ path slots while noFiLM-Tf uses only $7.9$. This indicates that the FiLM policy can make finer-grained, more differentiated path selections based on current link states and each flow's intent, whereas noFiLM-Tf, unable to perceive intent, can only rely on state and concentrates its selection on fewer path slots. This comparison corroborates the sensitivity evidence from the representation level: \emph{intent injected into representations by FiLM drives more diverse path selection}, linking representation-level differentiation to decision-level behavior.

